\section{Derivation of the radial functions}
\label{app:radial}

In this appendix we derive the integral representation \eqref{eq:Gab} of the functions
$\cG^{(\alpha\beta)}$ and the closed form \eqref{eq:GE} of $\cG_E$. Using eq.~\eqref{eq:notation}, the
product of the two energy denominators of eq.~\eqref{eq:Gdef} is written with two Schwinger parameters,
\begin{align}
&\frac{1}{\big(\kt^2+\ef^2\big)\Big[\zeta(1-\zeta)^2\big(\pt^{\,2}+\of^2\big)+\xi\eta\big(\kt^2+\ef^2\big)\Big]}
\nonumber\\
&\hspace{3cm}=\frac{1}{\xi\eta}\int_0^\infty\dd s_1\int_0^\infty\dd s_2\;
e^{-(s_1+s_2)(\kt^2+\ef^2)-\frac{s_2}{a}(\pt^{\,2}+\of^2)}\,.
\end{align}
The Gaussian integrals give
\begin{equation}
\int\frac{\dd^2\kt}{2\pi}\,e^{i\kt\cdot\R-(s_1+s_2)\kt^2}=\frac{e^{-\R^2/4(s_1+s_2)}}{2(s_1+s_2)}\,,\qquad
\int\frac{\dd^2\pt}{2\pi}\,e^{-i\pt\cdot\Y-\frac{s_2}{a}\pt^{\,2}}=\frac{a}{2s_2}\,e^{-a\Y^2/4s_2}\,.
\end{equation}
We change variables to the total Schwinger time $s=s_1+s_2$ and to the fraction $u=s_2/s\in[0,1]$ spent
in the three-parton state, with $\dd s_1\dd s_2=s\,\dd s\,\dd u$. The scalar integral becomes
\begin{align}
I(\R,\Y)&\equiv\int\frac{\dd^2\kt}{2\pi}\frac{\dd^2\pt}{2\pi}\;
\frac{e^{i\kt\cdot\R-i\pt\cdot\Y}}{\big(\kt^2+\ef^2\big)\Big[\zeta(1-\zeta)^2\big(\pt^{\,2}+\of^2\big)+\xi\eta\big(\kt^2+\ef^2\big)\Big]}
\nonumber\\
&=\frac{a}{4\xi\eta}\int_0^1\frac{\dd u}{u}\int_0^\infty\frac{\dd s}{s}\;
\exp\left\{-s\,A_u-\frac{\Delta_u}{4s}\right\}
=\frac{a}{2\xi\eta}\int_0^1\frac{\dd u}{u}\;K_0\Big(\sqrt{A_u\Delta_u}\Big)\,,
\label{eq:I}
\end{align}
with $A_u=\ef^2+u\,\of^2/a$ and $\Delta_u=\R^2+a\Y^2/u$ as in eq.~\eqref{eq:ADelta}, and where we used
$\int_0^\infty\frac{\dd s}{s}\,e^{-sA-\Delta/4s}=2K_0(\sqrt{A\Delta})$. Since
$(1-\zeta)\,a/(2\xi\eta)=1/[2\zeta(1-\zeta)]$, this is $\cG^{(00)}=(1-\zeta)I$. The other three functions
follow by differentiation, $\kt^j\to-i\partial/\partial R^j$ and $\pt^{\,n}\to i\partial/\partial Y^n$,
with $\partial\Delta_u/\partial R^j=2R^j$, $\partial\Delta_u/\partial Y^n=2aY^n/u$ and
\begin{equation}
\frac{\partial}{\partial\Delta}K_0\big(\sqrt{A\Delta}\big)=-\frac12\sqrt{\frac{A}{\Delta}}\,K_1\big(\sqrt{A\Delta}\big)\,,\qquad
\frac{\partial}{\partial\Delta}\left[\sqrt{\frac{A}{\Delta}}\,K_1\big(\sqrt{A\Delta}\big)\right]
=-\frac12\,\frac{A}{\Delta}\,K_2\big(\sqrt{A\Delta}\big)\,.
\end{equation}
This gives eq.~\eqref{eq:Gdef} with the functions \eqref{eq:Gexplicit}. For the instantaneous term
there is a single denominator, hence a single Schwinger parameter,
\begin{align}
\int\frac{\dd^2\kt}{2\pi}\frac{\dd^2\pt}{2\pi}\;&
\frac{e^{i\kt\cdot\R-i\pt\cdot\Y}}{\zeta(1-\zeta)^2\big(\pt^{\,2}+\of^2\big)+\xi\eta\big(\kt^2+\ef^2\big)}
\nonumber\\
&=\frac{a}{4\xi\eta}\int_0^\infty\frac{\dd s}{s^2}\;
\exp\left\{-s\Big(\ef^2+\frac{\of^2}{a}\Big)-\frac{\R^2+a\Y^2}{4s}\right\}
\nonumber\\
&=\frac{1}{1-\zeta}\;\frac{\Ef\,K_1(\Ef\,\cD)}{\cD}\,,
\end{align}
where we used $\int_0^\infty\frac{\dd s}{s^2}\,e^{-sA-\Delta/4s}=4\sqrt{A/\Delta}\,K_1(\sqrt{A\Delta})$,
together with $\ef^2+\of^2/a=\zeta(1-\zeta)\Ef^2$, which is eq.~\eqref{eq:Ef}, and
$\R^2+a\Y^2=\cD^2/[\zeta(1-\zeta)]$.

For a numerical evaluation of eq.~\eqref{eq:Gab} it is convenient to use the variable $t=-\ln u$: the
integrand is then a smooth function of $t$ which is peaked around $t\simeq\ln[\R^2/(a\Y^2)]$ when
$a\Y^2\ll\R^2$ and decays exponentially on both sides.

\section{Summary of the checks}
\label{app:checks}

All the formulas of this paper have been verified independently of their derivation, with computer
algebra and with numerical integration. We list here the checks that were performed.
\begin{enumerate}
\item \emph{Vertices and amplitude.} The complete tree-level amplitude for $\gT\to\qqg$ before the
collision, including the instantaneous pieces of the fermion propagators, was computed from Feynman
numerators with four-component Lepage--Brodsky spinors~\cite{Lepage:1980fj} of mass $\mf$, for all the
sixteen combinations of helicities and polarisations, for several values of $(\zeta,\xi,Q^2,\mf)$
including $\mf=0$, and for random transverse momenta, using exact rational arithmetic. It agrees with
the two-component expressions of section~\ref{sec:wf}, separately for each emitter and for their sum,
with a single overall constant. This fixes the helicity-flip terms of eqs.~\eqref{eq:Hgamma2} and
\eqref{eq:Hg2}, the instantaneous terms, and the relative sign and normalisation of the two emitters.
\item \emph{Radial functions.} The representation \eqref{eq:Gab} was compared with a direct numerical
evaluation of the Fourier transform \eqref{eq:Gdef}, in which the $\pt$ integral is done analytically
and the $\kt$ integral is a one-dimensional Hankel transform; the two agree to better than $10^{-10}$
for all four functions. Equation~\eqref{eq:GE} and the short-distance limits \eqref{eq:UV} were checked
in the same way.
\item \emph{Helicity sums.} Equations~\eqref{eq:H-direct}, \eqref{eq:Ht-direct} and \eqref{eq:Hx} were
obtained by summing the products of the effective vertices over the four helicity states and the four
polarisation states, symbolically in $\zeta$, $\xi$ and $\mf$, and by decomposing the result on a basis
of rotationally invariant tensors. The momentum-space sums \eqref{eq:H-el-qbar}, \eqref{eq:H-el-q} and
\eqref{eq:H-el-int} were checked against the amplitudes \eqref{eq:amp} for random complex values of the
transverse structures, with exact arithmetic.
\item \emph{Kinematics.} The changes of variables \eqref{eq:var-qbar}--\eqref{eq:relmom}, the shifted
momenta \eqref{eq:shifted} and the identities \eqref{eq:Ef} and \eqref{eq:f12} were verified
symbolically.
\item \emph{Limits.} The massless limit \eqref{eq:massless-fns}, the soft limit \eqref{eq:G-soft} and
the finiteness of the photoproduction limit were checked numerically; the common soft weights
\eqref{eq:soft-weights} and the collinear factorisation \eqref{eq:coll-H}--\eqref{eq:Pmassive} were
checked symbolically.
\item \emph{Normalisation.} The overall normalisation of the cross-section agrees with
ref.~\cite{Iancu:2022gpw} for $\mf=0$, and it reproduces the standard soft-gluon spectrum and the
standard quasi-collinear emission probability \eqref{eq:dN} in the two limits of
section~\ref{sec:limits}.
\end{enumerate}
